\documentclass[twocolumn]{article}
\usepackage{amsmath,amssymb,amsthm}

\makeatletter
% Plain marker as in paper-plan.md, plus a mode-tagged twin (spike addition:
% \ifvmode tells the extractor whether the destination was emitted in
% vertical mode -- between blocks -- or inside a line).
\newcommand\laxmark[2]{%
  \pdfdest name{lax.#2.#1} xyz\relax
  \ifvmode\pdfdest name{lax.#2.#1.v} xyz\relax
  \else\pdfdest name{lax.#2.#1.h} xyz\relax\fi}
\makeatother

\newtheorem{definition}{Definition}
\newtheorem{theorem}{Theorem}

\newcommand{\fillerA}{The study of structural graph parameters has a long
history, and the notion of tree decompositions in particular occupies a
central place in it. A graph that admits a decomposition whose bags are
uniformly small behaves, for many algorithmic purposes, almost like a tree,
and a great number of problems that are intractable in general become
solvable in linear time on such inputs. We recall the basic notions here,
fix the notation used throughout, and state the results that the remainder
of this note develops in detail.}

\newcommand{\fillerB}{Throughout, all graphs are finite, simple and
undirected. For a graph $G$ we write $V(G)$ for its vertex set and $E(G)$
for its edge set, and we abbreviate $|V(G)|$ by $n$ whenever the graph is
clear from the context. A separation of $G$ is a pair of vertex sets whose
union is $V(G)$ and between which there is no edge outside their
intersection; the order of a separation is the size of that intersection.
Separations of small order are the basic building blocks of every
decomposition considered below.}

\newcommand{\fillerC}{It is worth emphasising that none of the arguments in
this section depend on the particular representation chosen for the
decomposition. Any of the usual encodings, whether by a tree with labelled
nodes, by a nested family of separations, or by a linear order on the
vertices with a bounded window, can be translated into any other with only
a constant factor loss in the width. We fix one of them for concreteness
and note in passing where the translation would matter.}

\begin{document}

\title{A Spike Paper on Decompositions}
\author{Lax Viewer Spike}
\date{}
\maketitle

\section{Introduction}

\fillerA

\fillerB

\laxmark{b}{1}%
\begin{definition}[Tree decomposition]
A \emph{tree decomposition} of a graph $G$ is a pair $(T,\mathcal{B})$
consisting of a tree $T$ and a family
\laxmark{b}{5}%
$\mathcal{B} = (B_t)_{t \in V(T)}$ of subsets of $V(G)$, called the bags of
the decomposition, such that every vertex and every edge is covered and the
bags containing any fixed vertex form a connected subtree of $T$\laxmark{e}{5}%
.
The width of the decomposition is the largest bag size minus one.
\end{definition}
\laxmark{e}{1}%

\fillerC

We now turn to the parameter itself. The minimum width of \laxmark{b}{2}%
a tree decomposition \laxmark{e}{2}%
of a graph is one of the most
thoroughly studied quantities in structural graph theory, and it is the
quantity that the whole of this note is concerned with.

\fillerA

\section{A range that crosses a page}

\laxmark{b}{3}%
\fillerB
\fillerC
\laxmark{e}{3}%

\fillerA

\section{A range that crosses a column}

\fillerB

\begin{theorem}
Every graph that admits a decomposition of width at most $k$ also admits
one of width at most $k$ in which the underlying tree has at most $n$
nodes, and such a decomposition can be computed in linear time from the
given one.
\end{theorem}

\fillerC

\laxmark{b}{4}%
\fillerA
\fillerB
\laxmark{e}{4}%

\fillerC

\fillerA

\section{Closing remarks}

The next range starts at a forced line break, so that its begin marker is
the first thing on its line and its destination lands on the column's left
edge -- exactly where a vertical-mode marker would land too.\\
\laxmark{b}{6}%
This whole line and the next few words belong to the passage, which ends
right here\laxmark{e}{6}%
, and the sentence carries on afterwards without interruption.

\fillerB

\fillerC

\fillerA

\end{document}
